\section{Discussion}
\model{} advances pegRNA ranking through edit-aware sequence modeling and order-sensitive supervision. Relative to released OptiPrime, its source ensemble improves pooled Spearman correlation, correlation within experimental conditions, correlation within protospacer groups and Kendall's tau-b. A positive-outcome-only analysis retains the advantage, and the ordinal--S4D ensemble extends higher pooled correlation to an external library. Together, these results show that the source ranking gain is not confined to a single metric, one source-study subset or the separation of zero from nonzero outcomes.

The architecture provides a coherent way to represent the design problem. Aligned wild-type/edited pairs expose the intended change, segment-aware pegRNA encoding distinguishes functional sequence regions, and cross-attention integrates guide and target information. Context conditioning supports a heterogeneous training corpus. Ordinal supervision supplies multiple outcome-threshold comparisons rather than requiring the ranking score itself to be a calibrated efficiency. The development comparisons support ordinal supervision and S4D as useful choices within this framework. They do not isolate a uniquely causal component: the historical comparisons are not a fully matched randomized factorial, and the systems also differ in training weights and ensemble composition. We accordingly attribute the benchmark improvement to the evaluated \model{} system, not exclusively to one architectural mechanism.

Importantly, the contribution extends beyond correlation. On informative source decisions, \model{} increases the measured efficiency of the selected pegRNA by 7.2\% relative to OptiPrime and raises the measured-best-design hit rate from 65.35\% to 74.23\%. Target-library fine-tuning then establishes an external selection advantage: native-ordinal adaptation improves mean top-choice efficiency by 3.9\% relative to released OptiPrime, with positive paired intervals in all three optimizer seeds. Ordinary ordinal adaptation is sufficient for this result; separate utility or multitask heads do not establish a replicated additional advantage over it. These findings favor a simple supported adaptation procedure over treating an unvalidated extra branch as an essential part of the method.

The external results also define when this selection claim applies. The source-trained ordinal--S4D ensemble has higher external pooled correlation but lower zero-shot top-choice efficiency than OptiPrime. Ranking across a library and ordering alternatives within a fixed edit emphasize different variations, particularly when candidate depth, sequence geometry and outcome distributions change. Fine-tuning demonstrates that target supervision can improve the relevant choices, but does not by itself identify which aspect of domain shift caused the zero-shot shortfall. Nor does it establish that ranking-specific losses are necessary: accurate conditional-mean estimates would also suffice for expected-efficiency selection. Our evidence supports a source-trained ranking advantage and a target-supervised selection advantage, not unconditional superiority on every deployment endpoint.

Label availability and source retention remain practical constraints. In the budget-accounted follow-up, adequately tuned ordinal adaptation and a source-anchored selector did not establish better external selection than OptiPrime at a nominal 1,000-group budget. Balanced teacher-margin preservation improved source retention relative to unregularized anchoring, but source-constrained checkpoint selection incurred a target-efficiency cost. These supplementary results concern additional low-label and retention requirements; they do not negate the full-budget adaptation result, nor justify extending it to lower budgets. Auditing the score actually used to select candidates was essential, because an unchanged prediction head could conceal forgetting in a separate deployed selector.

Several qualifications delimit the comparisons. The source benchmark was revisited during model development and subsequent audits. The external library was initially reserved, but later reused extensively; even the full-budget adaptation test is a retrospective within-library evaluation rather than a new independent study. Moreover, OptiPrime did not receive the target labels used for \model{} fine-tuning. The external selection win is therefore a practical comparison with a released tool, not an equal-label comparison of adaptation methods. The source selection population is shallow and Kim-derived, while pooled ranking also includes Liu-derived records. Repeated optimizer seeds quantify one form of training variability, not generalization across laboratories or studies. Measured maxima contain assay noise and are not noise-free biological optima. Our claims concern the stated populations and released comparator, not every application of OptiPrime \citep{hsu2026}.

Reproducible deployment also requires explicit model identity. The heterogeneous source ensemble, five-fold ordinal--S4D transfer ensemble and fine-tuned single backbone are different configurations of the framework. Population-rank averaging in the historical source ensemble requires a fixed reference transformation or a validated alternative interface for isolated queries. Legacy RNA/DNA tokenization must be handled consistently. The harmonized corpus includes author-supplied material with unresolved redistribution permissions, and the recorded inference timings do not establish a speed advantage over OptiPrime. The accompanying aggregate evidence, buildable manuscript and deployment audits document these boundaries without substituting for a complete public-data release.

The next computational confirmation should freeze the supported ordinal adaptation procedure, account for both training and validation labels, and evaluate ranking and deployed top-choice efficiency on verified independent fixed-edit decisions. An equal-target-label comparator and a source-training study holdout would address fairness and cross-study transfer separately. Higher-label learning curves and carefully controlled representation changes could then determine how far the external selection advantage extends. These are future tests, not completed evidence.

In conclusion, \model{} is a new pegRNA-ranking framework that outperforms released OptiPrime across the evaluated source rank-based metrics and retains higher pooled ranking accuracy on an external library. Its ranking gains translate into improved source-domain guide selection, and native-ordinal fine-tuning enables an additional external-library selection-efficiency advantage. These are the central contributions of the method; the accompanying controls establish the conditions under which each is supported.
